%************************************************
\chapter{Complex Numbers from Scratch}\label{app:complex}
%************************************************

This appendix is a short course on the complex numbers, for readers who would like to see the rules of Chapter~\ref{ch:one-number} in one place. It follows the teaching approach of the lecture notes of \citet{MokhithiShock2020}, released under the Creative Commons Attribution 4.0 license; the text here is a new and shorter presentation. For a beautiful, picture-first treatment, see \citet{Needham1997}.

%------------------------------------------------
\section{A number whose square is minus one}\label{app:c-i}
%------------------------------------------------

Every number on the number line gives a result of zero or more when it is multiplied by itself. So the equation $x^2 = -1$ has no solution on the line. Instead of giving up, we \emph{define} a new number.

\begin{definition}[The imaginary unit]\label{def:imaginary-unit}
The imaginary unit $i$ is a number whose square is $-1$:
\[
i^2 = -1 .
\]
Its opposite, $-i$, has the same property: $(-i)^2 = -1$. So $i$ and $-i$ are the two numbers which, when squared, give $-1$.
\end{definition}

A \emph{complex number} is a number of the form
\[
z = a + bi, \qquad a, b \in \R .
\]
The number $a$ is the \emph{real part} of $z$, written $\operatorname{Re} z$, and $b$ is its \emph{imaginary part}, written $\operatorname{Im} z$. Two complex numbers are equal when their real parts are equal and their imaginary parts are equal. The real numbers are the complex numbers with $b = 0$, so nothing old is lost.

%------------------------------------------------
\section{Arithmetic}\label{app:c-arithmetic}
%------------------------------------------------

Complex numbers are added and multiplied with the ordinary rules of algebra, remembering that $i^2 = -1$:
\begin{align*}
(a + bi) + (c + di) &= (a + c) + (b + d)\,i, \\
(a + bi)(c + di) &= ac + adi + bci + bd\,i^2 \\
&= (ac - bd) + (ad + bc)\,i .
\end{align*}
For example, $(2 + i)^2 = 4 + 4i + i^2 = 3 + 4i$, and $(2 + i)^3 = (3 + 4i)(2 + i) = 6 + 3i + 8i + 4i^2 = 2 + 11i$. This is the calculation with which Bombelli tamed his cubic equation in 1572 (Chapter~\ref{ch:one-number}).

The \emph{conjugate} of $z = a + bi$ is $\bar{z} = a - bi$. Multiplying a number by its conjugate always gives a real number that is zero or more:
\[
z\bar{z} = (a + bi)(a - bi) = a^2 + b^2 .
\]
This is the key to division. To divide by $c + di$, multiply the top and the bottom by its conjugate $c - di$:
\[
\frac{a + bi}{c + di} = \frac{(a + bi)(c - di)}{(c + di)(c - di)} = \frac{(ac + bd) + (bc - ad)\,i}{c^2 + d^2} .
\]

%------------------------------------------------
\section{The plane}\label{app:c-plane}
%------------------------------------------------

A complex number $a + bi$ is a point of the plane: $a$ steps along the real axis and $b$ steps along the imaginary axis (\Cref{fig:complex-plane}). Two numbers describe the point in another way.
\begin{itemize}
  \item The \emph{modulus} $|z| = \sqrt{a^2 + b^2}$ is the distance from $0$ to $z$, by the theorem of Pythagoras. Note that $|z|^2 = z\bar{z}$.
  \item The \emph{argument} $\theta$ is the angle from the positive real axis to the segment from $0$ to $z$, measured counterclockwise. It satisfies $\tan\theta = b/a$, and the signs of $a$ and $b$ tell in which quarter of the plane the angle lies.
\end{itemize}
With $r = |z|$, the number is written in \emph{polar form}:
\[
z = r(\cos\theta + i\sin\theta) .
\]
The conjugate $\bar{z}$ is the mirror image of $z$ in the real axis: same modulus, opposite argument.

\begin{figure}[tb]
\myfloatalign
\begin{tikzpicture}[>=Stealth, scale=1.1, every node/.style={font=\small}]
  \draw[->, gray] (-0.6,0) -- (4.4,0) node[right] {real};
  \draw[->, gray] (0,-2.4) -- (0,2.7) node[above] {imaginary};
  \coordinate (Z) at (3,2);
  \coordinate (Zb) at (3,-2);
  \draw[gray, dashed] (Z) -- (3,0) node[below right] {$a$};
  \draw[gray, dashed] (Z) -- (0,2) node[left] {$b$};
  \draw[very thick, Maroon] (0,0) -- (Z) node[midway, above left] {$r = |z|$};
  \draw[thick, RoyalBlue, dashed] (0,0) -- (Zb);
  \draw (0.8,0) arc (0:33.69:0.8);
  \node at (16.8:1.1) {$\theta$};
  \fill[Maroon] (Z) circle (1.6pt) node[above right] {$z = a + bi$};
  \fill[RoyalBlue] (Zb) circle (1.6pt) node[below right] {$\bar{z} = a - bi$};
  \fill (0,0) circle (1.2pt) node[below left] {$0$};
\end{tikzpicture}
\caption[A complex number in the plane]{A complex number $z = a + bi$ in the plane, with its modulus $r = |z|$, its argument $\theta$, and its conjugate $\bar{z}$, the mirror image of $z$ in the real axis.}
\label{fig:complex-plane}
\end{figure}

%------------------------------------------------
\section{Multiplying means turning}\label{app:c-turning}
%------------------------------------------------

Take two numbers in polar form, $z_1 = r_1(\cos\theta_1 + i\sin\theta_1)$ and $z_2 = r_2(\cos\theta_2 + i\sin\theta_2)$. Multiplying them out and using the addition formulas of trigonometry gives
\[
z_1 z_2 = r_1 r_2 \bigl(\cos(\theta_1 + \theta_2) + i\sin(\theta_1 + \theta_2)\bigr).
\]
This is the rule that Wessel found in 1797: \emph{to multiply two complex numbers, multiply their lengths and add their angles}. In particular, $i = \cos(\pi/2) + i\sin(\pi/2)$ has length one and angle $\pi/2$, so multiplying by $i$ turns any number by a quarter turn, and multiplying by $i$ twice turns it by a half turn, which is multiplication by $-1$. That is the whole meaning of $i^2 = -1$.

%------------------------------------------------
\section{Euler's formula and de Moivre's formula}\label{app:c-euler}
%------------------------------------------------

Euler discovered that the exponential function and the circle are the same story:
\[
e^{i\theta} = \cos\theta + i\sin\theta .
\]
Every complex number can therefore be written as $z = re^{i\theta}$, and the multiplication rule becomes the ordinary rule for exponents: $r_1e^{i\theta_1} \cdot r_2e^{i\theta_2} = r_1r_2\,e^{i(\theta_1 + \theta_2)}$. Raising a number of length one to the power $n$ turns it $n$ times, which is \emph{de Moivre's formula}:
\[
(\cos\theta + i\sin\theta)^n = \cos(n\theta) + i\sin(n\theta) .
\]
With $\theta = \pi$, Euler's formula gives $e^{i\pi} = -1$, often written $e^{i\pi} + 1 = 0$: one short line that joins five of the most important numbers of mathematics.

%------------------------------------------------
\section{Roots of unity}\label{app:c-roots}
%------------------------------------------------

The equation $z^n = 1$ has exactly $n$ complex solutions:
\[
z_k = e^{2\pi i k/n} = \cos\frac{2\pi k}{n} + i\sin\frac{2\pi k}{n}, \qquad k = 0, 1, \dots, n - 1 .
\]
They are the corners of a regular polygon with $n$ sides, drawn inside the circle of radius one with one corner at $1$. For $n = 4$ they are $1$, $i$, $-1$ and $-i$: the cycle of powers of $i$ from Chapter~\ref{ch:one-number}. In Chapter~\ref{ch:general-numbers}, the powers of the Composer Unit at $\pb = \pi/3$ turn by the same steps as the sixth roots of unity, but grow at every step (\Cref{fig:spiral}).

%------------------------------------------------
\section{Exercises}\label{app:c-exercises}
%------------------------------------------------

\begin{enumerate}
  \item Compute $(1 + i)^2$ and $(1 + i)^4$.
  \item Find the modulus and the argument of $-1 + i\sqrt{3}$.
  \item Compute $\dfrac{3 + 4i}{1 + 2i}$.
  \item Using polar form, explain why $i^3 = -i$ and $i^4 = 1$.
\end{enumerate}

\begin{workedexample}[Answers]
\begin{enumerate}
  \item $(1 + i)^2 = 1 + 2i + i^2 = 2i$, and $(1 + i)^4 = (2i)^2 = -4$.
  \item The modulus is $\sqrt{1 + 3} = 2$. Since $\cos\theta = -1/2$ and $\sin\theta = \sqrt{3}/2$, the argument is $\theta = 2\pi/3$.
  \item $\dfrac{(3 + 4i)(1 - 2i)}{(1 + 2i)(1 - 2i)} = \dfrac{3 - 6i + 4i - 8i^2}{5} = \dfrac{11 - 2i}{5}$.
  \item The number $i$ has length one and angle $\pi/2$. Its third power has angle $3\pi/2$, which points straight down: $-i$. Its fourth power has angle $2\pi$, a full turn: $1$.
\end{enumerate}
\end{workedexample}
